% Partie : théorie de la démonstration
% Chapitre : calcul des séquents
% Section : inversibilité

\documentclass[../../../include/open-logic-section]{subfiles}

\begin{document}

\olfileid[fr]{pt}{seq}{inv}

\olsection{Inversibilité des règles}

\begin{defn}
Une règle~$R$,
\[
\AxiomC{$S_1$}
\AxiomC{$\dots$}
\AxiomC{$S_n$}
\RightLabel{$R$}
\TrinaryInfC{$S$}
\DisplayProof
\]
est \emph{inversible} dans un calcul si, dès que
$\Proves S$, on a $\Proves S_i$ pour
$i = 1$, \dots,~$n$.
Elle est \emph{inversible sans augmentation de hauteur},
ou \emph{hp-inversible}, si, pour toute borne de hauteur,
$\Proves[h] S$ entraîne $\Proves[h] S_i$
pour $i = 1$, \dots,~$n$.
La borne de hauteur et le nombre de prémisses
sont ici deux paramètres distincts.
\end{defn}

\begin{lem}\ollabel{lem:G3c-invert}
Les règles de~\Log{G3c} pour $\lnot$, $\land$,
$\lor$ et $\lif$ sont inversibles sans augmentation
de hauteur. Plus précisément :
\begin{enumerate}
  \item \RightR{\lnot} : si
  $\Log{G3c} \Proves[n] \Gamma \Sequent \Delta, \lnot !A$,
  alors $\Log{G3c} \Proves[n] !A, \Gamma \Sequent \Delta$.
  \item \LeftR{\lnot} : si
  $\Log{G3c} \Proves[n] \Gamma, \lnot !A \Sequent \Delta$,
  alors $\Log{G3c} \Proves[n] \Gamma \Sequent \Delta, !A$.
  \item \RightR{\land} : si
  $\Log{G3c} \Proves[n] \Gamma \Sequent \Delta, !A \land !B$,
  alors $\Log{G3c} \Proves[n] \Gamma \Sequent \Delta, !A$
  et $\Log{G3c} \Proves[n] \Gamma \Sequent \Delta, !B$.
  \item \LeftR{\land} : si
  $\Log{G3c} \Proves[n] \Gamma, !A \land !B \Sequent \Delta$,
  alors $\Log{G3c} \Proves[n] !A, !B, \Gamma \Sequent \Delta$.
  \item \RightR{\lor} : si
  $\Log{G3c} \Proves[n] \Gamma \Sequent \Delta, !A \lor !B$,
  alors $\Log{G3c} \Proves[n] \Gamma \Sequent \Delta, !A, !B$.
  \item \LeftR{\lor} : si
  $\Log{G3c} \Proves[n] !A \lor !B, \Gamma \Sequent \Delta$,
  alors $\Log{G3c} \Proves[n] !A, \Gamma \Sequent \Delta$
  et $\Log{G3c} \Proves[n] !B, \Gamma \Sequent \Delta$.
  \item \RightR{\lif} : si
  $\Log{G3c} \Proves[n] \Gamma \Sequent \Delta, !A \lif !B$,
  alors $\Log{G3c} \Proves[n] !A, \Gamma \Sequent \Delta, !B$.
  \item \LeftR{\lif} : si
  $\Log{G3c} \Proves[n] !A \lif !B, \Gamma \Sequent \Delta$,
  alors $\Log{G3c} \Proves[n] \Gamma \Sequent \Delta, !A$
  et $\Log{G3c} \Proves[n] !B, \Gamma \Sequent \Delta$.
\end{enumerate}
\end{lem}

\begin{proof}
Pour chacune de ces règles~$R$, une preuve~$\pi$
de sa conclusion~$S$ doit fournir des preuves~$\pi_i$
de ses prémisses $S_i$, avec
$\pheight{\pi_i} \le \pheight{\pi}$.
Considérons \LeftR{\land}. Supposons donnée une
preuve~$\pi$ de $!A \land !B, \Gamma \Sequent \Delta$.
Il faut construire une preuve $\pi'$ de
$!A, !B, \Gamma \Sequent \Delta$ telle que
$\pheight{\pi'} \le \pheight{\pi}$.
On raisonne par récurrence sur $n = \pheight{\pi}$.

Si $n = 0$, la preuve $\pi$ se réduit à l'axiome
$!A \land !B, \Gamma \Sequent \Delta$.
Un axiome d'identité de \Log{G3c} a la forme
$!C, \Gamma \Sequent \Delta, !C$, avec $!C$ atomique.
La formule $!C$ ne peut donc pas être $!A \land !B$ :
elle se trouve à la fois dans $\Gamma$ et~$\Delta$.
Le séquent $!A, !B, \Gamma \Sequent \Delta$
reste alors un axiome d'identité.
Si l'axiome initial est un axiome de l'absurde,
l'absurde se trouve dans l'antécédent de contexte,
et y reste après la décomposition.
Dans les deux cas, on obtient une preuve~$\pi'$
de hauteur~$0$.

Si $n>0$, la preuve se termine par une inférence.
L'hypothèse de récurrence vaut pour toute preuve
de hauteur~$<n$, quels que soient ses contextes :
pour $k<n$ et tous $\Pi$ et $\Lambda$,
une preuve de hauteur~$k$ de
$!A \land !B, \Pi \Sequent \Lambda$ donne une
preuve de hauteur~$\le k$ de
$!A, !B, \Pi \Sequent \Lambda$.

Si l'occurrence choisie de $!A \land !B$ est
principale dans la dernière inférence, celle-ci
est \LeftR{\land}. Sa sous-preuve immédiate~$\pi_1$
se termine par $!A, !B, \Gamma \Sequent \Delta$,
et $\pheight{\pi_1} < \pheight{\pi}$ donne le résultat.

Sinon, $!A \land !B$ appartient au contexte.
On applique l'hypothèse aux sous-preuves immédiates
de la dernière règle~$R$. On obtient une preuve~$\pi_1'$,
ou deux preuves~$\pi_1'$ et~$\pi_2'$,
de hauteurs au plus celles des sous-preuves originales.
Leurs séquents finaux ont $!A, !B$ à gauche
à la place de $!A \land !B$. On réapplique $R$
pour obtenir~$\pi'$. Les constantes de ces formules
sont celles de la conjonction : cette décomposition
n'en introduit aucune et conserve les conditions
sur les constantes propres.
Par exemple, si $\pi$ se termine par~\LeftR{\lif},
\[
\AxiomC{}
\RightLabel{$\pi_1$}
\Deduce$!A \land !B, \Gamma' \fCenter \Delta, !C$
\AxiomC{}
\RightLabel{$\pi_2$}
\Deduce$!A \land !B, \Gamma', !D \fCenter \Delta$
\RightLabel{\LeftR{\lif}}
\BinaryInf$!A \land !B, \Gamma', !C \lif !D \fCenter \Delta$
\DisplayProof
\]
la principale est $!C \lif !D$, et le contexte
gauche est $!A \land !B, \Gamma'$,
avec $\Gamma = \Gamma', !C \lif !D$.
L'hypothèse donne des preuves $\pi_1'$ et $\pi_2'$
de $!A, !B, \Gamma' \Sequent \Delta, !C$
et de $!A, !B, \Gamma', !D \Sequent \Delta$.
On les combine par \LeftR{\lif} :
\[
\AxiomC{}
\RightLabel{$\pi_1'$}
\Deduce$!A, !B, \Gamma' \fCenter \Delta, !C$
\AxiomC{}
\RightLabel{$\pi_2'$}
\Deduce$!A, !B, \Gamma', !D \fCenter \Delta$
\RightLabel{\LeftR{\lif}}
\BinaryInf$!A, !B, \Gamma', !C \lif !D \fCenter \Delta$
\DisplayProof
\]
Les hauteurs vérifient
\begin{align*}
\pheight{\pi} & = \max(\pheight{\pi_1}, \pheight{\pi_2}) + 1\\
\pheight{\pi'} & = \max(\pheight{\pi_1'}, \pheight{\pi_2'}) + 1
\intertext{par définition, et}
\pheight{\pi_1'} & \le \pheight{\pi_1}\\
\pheight{\pi_2'} & \le \pheight{\pi_2}
\intertext{par hypothèse de récurrence. Donc,}
\pheight{\pi'} & \le \pheight{\pi}.
\end{align*}
\end{proof}

\begin{prob}
Démontrer le \olref{lem:G3c-invert} pour \RightR{\lif}.
\end{prob}

\begin{lem}\ollabel{lem:invert-quant}
Les règles \RightR{\lforall} et \LeftR{\lexists}
sont inversibles sans augmentation de hauteur
au sens suivant :
\begin{enumerate}
  \item si
  $\Log{G3c} \Proves[n] \Gamma \Sequent \Delta, \lforall[x][!A(x)]$,
  alors $\Log{G3c} \Proves[n] \Gamma \Sequent \Delta, !A(c)$ ;
  \item si
  $\Log{G3c} \Proves[n] \lexists[x][!A(x)], \Gamma \Sequent \Delta$,
  alors $\Log{G3c} \Proves[n] !A(c), \Gamma \Sequent \Delta$ ;
\end{enumerate}
pour toute constante $c$ absente de $\Gamma$,
de $\Delta$ et, respectivement, de
$\lforall[x][!A(x)]$ ou de $\lexists[x][!A(x)]$.
\end{lem}

\begin{proof}
On démontre (1). Le cas (2) est laissé en exercice.
On peut supposer la preuve~$\pi$ régulière et choisir une
constante~$d$, distincte de~$c$, qui n'apparaît nulle part
dans~$\pi$, y compris comme constante propre.
On obtient d'abord une preuve de $\Gamma \Sequent \Delta, !A(d)$,
par récurrence sur la hauteur de~$\pi$ comme dans le
\olref{lem:G3c-invert}.

À hauteur zéro, la formule quantifiée n'est pas atomique.
L'identité atomique qui ferme le séquent reste dans les contextes,
ou l'absurde reste à gauche. L'instance obtenue est donc un axiome.

Si $\lforall[x][!A(x)]$ est principale, la dernière règle est
$\RightR{\lforall}$ et sa prémisse est
$\Gamma \Sequent \Delta, !A(a)$. La sous-preuve immédiate~$\pi_1$
a une hauteur strictement inférieure à celle de~$\pi$.
La constante~$a$ est absente de la conclusion de l'inférence.
Comme $\pi$ est régulière, $a$ n'est constante propre d'aucune
inférence de~$\pi_1$. La constante~$d$ n'en est pas une non plus.
Le \olref[qua]{lem:subst-regular} s'applique à
$\Subst{\pi_1}{d}{a}$. On obtient une preuve de même hauteur
que~$\pi_1$ du séquent $\Gamma \Sequent \Delta, !A(d)$.

Si $\lforall[x][!A(x)]$ n'est pas principale, elle appartient
au contexte de la dernière inférence. Par exemple, pour
\RightR{\land} avec $!B \land !C$ principale, la preuve se termine par
\[
\AxiomC{}
\RightLabel{$\pi_1$}
\Deduce$\Gamma \fCenter \Delta', \lforall[x][!A(x)], !B$
\AxiomC{}
\RightLabel{$\pi_2$}
\Deduce$\Gamma \fCenter \Delta', \lforall[x][!A(x)], !C$
\RightLabel{\RightR{\land}}
\BinaryInf$\Gamma \fCenter \Delta', \lforall[x][!A(x)], !B \land !C$
\DisplayProof
\]
où $\Delta = \Delta', !B \land !C$.
Puisque $d$ est absente des deux sous-preuves, l'hypothèse
de récurrence fournit $\pi_1'$ et~$\pi_2'$ telles que
$\pheight{\pi_1'} \le \pheight{\pi_1}$ et
$\pheight{\pi_2'} \le \pheight{\pi_2}$. On obtient
\[
\AxiomC{}
\RightLabel{$\pi_1'$}
\Deduce$\Gamma \fCenter \Delta', !A(d), !B$
\AxiomC{}
\RightLabel{$\pi_2'$}
\Deduce$\Gamma \fCenter \Delta', !A(d), !C$
\RightLabel{\RightR{\land}}
\BinaryInf$\Gamma \fCenter \Delta', !A(d), !B \land !C$
\DisplayProof
\]
avec
$\pheight{\pi'} = \max(\pheight{\pi_1'}, \pheight{\pi_2'})+1 \le \pheight{\pi}$.
Les autres règles se réappliquent de la même façon. Pour
\LeftR{\lforall} et \RightR{\lexists}, on garde le même terme
d'instanciation, qui ne contient pas~$d$. Pour les règles
à constante propre, l'instance $!A(d)$ ne contient pas
la constante propre de la dernière inférence, puisque
celle-ci était absente de la conclusion et distincte de~$d$.

On a ainsi une preuve~$\pi'$ de $\Gamma \Sequent \Delta, !A(d)$
de hauteur au plus $\pheight{\pi}$. On la régularise en choisissant
toutes ses constantes propres différentes de $c$ et de~$d$,
sans changer sa conclusion ni sa hauteur. On applique ensuite
le lemme de substitution à $\Subst{\pi'}{c}{d}$. Comme $d$
était absente de la conclusion de~$\pi$, le séquent obtenu
est bien $\Gamma \Sequent \Delta, !A(c)$.
%
%
%
%
%
\end{proof}

\Olref{lem:G3c-invert} intervient dans la preuve
du théorème d'élimination des coupures.
Il permet aussi d'obtenir le résultat suivant.

\begin{prop}\ollabel{prop:G3c-cont-adm}
Les règles de contraction \LeftR{\Contraction}
et \RightR{\Contraction} de~\Log{G1c}
sont admissibles dans~\Log{G3c} sans augmentation
de hauteur.
\end{prop}

\begin{proof}
On traite simultanément \RightR{\Contraction}
et \LeftR{\Contraction}.
Soit $\pi$ une preuve de
$\Gamma \Sequent \Delta, !A, !A$ ou de
$!A, !A, \Gamma \Sequent \Delta$ dans~\Log{G3c}.
Il faut une preuve~$\pi'$ dans \Log{G3c} de
$\Gamma \Sequent \Delta, !A$ ou, respectivement,
de $!A, \Gamma \Sequent \Delta$, avec
$\pheight{\pi'} \le \pheight{\pi}$.
On raisonne par récurrence sur $n = \pheight{\pi}$.

Si $n=0$, la preuve se réduit à un axiome.
Lorsque $\Gamma \Sequent \Delta, !A, !A$
est un axiome d'identité, une formule atomique~$!C$
est commune à $\Gamma$ et~$\Delta$, ou $!A$
est atomique et appartient à~$\Gamma$.
Après suppression d'une des deux occurrences,
$\Gamma \Sequent \Delta, !A$ reste un axiome.
Un axiome de l'absurde reste aussi un axiome,
puisque son antécédent n'est pas modifié.
Pour $!A, !A, \Gamma \Sequent \Delta$,
le même argument conserve une identité atomique ;
si l'absurde est la formule contractée, une occurrence
reste à gauche. Les deux types d'axiomes sont donc couverts.

Si $n>0$, soit $R$ la dernière règle.
Quand les occurrences à contracter de $!A$
ne sont pas principales, on applique l'hypothèse
aux sous-preuves immédiates puis on réapplique $R$,
comme dans le \olref{lem:G3c-invert}.
L'hypothèse porte sur toute preuve de hauteur~$k<n$
de $\Pi \Sequent \Lambda, !D, !D$
ou de $!D, !D, \Pi \Sequent \Lambda$.
Elle donne une preuve de hauteur~$\le k$ de
$\Pi \Sequent \Lambda, !D$ ou, respectivement,
de $!D, \Pi \Sequent \Lambda$.
Les contextes et la formule contractée peuvent donc
différer de $\Gamma$, de $\Delta$ et de~$!A$.

Par exemple, si $R$ est $\RightR{\land}$
et que $\pi$ se termine par
\[
\AxiomC{}
\RightLabel{$\pi_1$}
\Deduce$\Gamma \fCenter \Delta', !B, !A, !A$
\AxiomC{}
\RightLabel{$\pi_2$}
\Deduce$\Gamma \fCenter \Delta', !C, !A, !A$
\RightLabel{\RightR{\land}}
\BinaryInf$\Gamma \fCenter \Delta', !B \land !C, !A, !A$
\DisplayProof
\]
avec $\Delta = \Delta', !B \land !C$,
l'hypothèse fournit des preuves $\pi_1'$ et $\pi_2'$
de $\Gamma \fCenter \Delta', !B, !A$
et de $\Gamma \fCenter \Delta', !C, !A$,
avec $\pheight{\pi_1'} \le \pheight{\pi_1}$
et $\pheight{\pi_2'} \le \pheight{\pi_2}$.
La preuve~$\pi'$ est
\[
\AxiomC{}
\RightLabel{$\pi_1'$}
\Deduce$\Gamma \fCenter \Delta', !B, !A$
\AxiomC{}
\RightLabel{$\pi_2'$}
\Deduce$\Gamma \fCenter \Delta', !C, !A$
\RightLabel{\RightR{\land}}
\BinaryInf$\Gamma \fCenter \Delta', !B \land !C, !A$
\DisplayProof
\]

Si une occurrence de $!A$ est principale,
on applique le lemme d'inversion
\olref{lem:G3c-invert} aux sous-preuves immédiates,
puis l'hypothèse de récurrence et enfin la règle~$R$.
Supposons, par exemple, que $\pi$ prouve
$\Gamma \Sequent \Delta, !A, !A$,
avec $!A \ident (!B \land !C)$ principale
dans la dernière inférence :
\[
\AxiomC{}
\RightLabel{$\pi_1$}
\Deduce$\Gamma \fCenter \Delta, !B \land !C, !B$
\AxiomC{}
\RightLabel{$\pi_2$}
\Deduce$\Gamma \fCenter \Delta, !B \land !C, !C$
\RightLabel{\RightR{\land}}
\BinaryInf$\Gamma \fCenter \Delta, !B \land !C, !B \land !C$
\DisplayProof
\]
Le \olref{lem:G3c-invert} appliqué à $\pi_1$
donne une preuve~$\pi_1'$ de
$\Gamma \Sequent \Delta, !B, !B$.
Il donne aussi une preuve de
$\Gamma \Sequent \Delta, !C, !B$,
dont nous n'avons pas besoin.
Appliqué à~$\pi_2$, le
\olref{lem:G3c-invert} fournit une preuve~$\pi_2'$
de $\Gamma \Sequent \Delta, !C, !C$.
L'inversion de \RightR{\land} n'augmente pas la hauteur :
$\pheight{\pi_1'} \le \pheight{\pi_1} < \pheight{\pi}$
et $\pheight{\pi_2'} \le \pheight{\pi_2} < \pheight{\pi}$.
L'hypothèse s'applique donc à $\pi_1'$ et $\pi_2'$.
On obtient des preuves~$\pi_1''$ et $\pi_2''$
de $\Gamma \Sequent \Delta, !B$
et de $\Gamma \Sequent \Delta, !C$, avec
$\pheight{\pi_1''} \le \pheight{\pi_1'} \le \pheight{\pi_1}$
et $\pheight{\pi_2''} \le \pheight{\pi_2'} \le \pheight{\pi_2}$.
La preuve~$\pi'$ est alors
\[
\AxiomC{}
\RightLabel{$\pi_1''$}
\Deduce$\Gamma \fCenter \Delta, !B$
\AxiomC{}
\RightLabel{$\pi_2''$}
\Deduce$\Gamma \fCenter \Delta, !C$
\RightLabel{\RightR{\land}}
\BinaryInf$\Gamma \fCenter \Delta, !B \land !C$
\DisplayProof
\]
avec $\pheight{\pi'} \le \pheight{\pi}$.

Les cas quantifiés où la formule contractée
est principale demandent une attention particulière.
Supposons $!A \ident \lforall[x][!B(x)]$
principale à droite. La preuve $\pi$ se termine par
\[
\AxiomC{}
\RightLabel{$\pi_1$}
\Deduce$\Gamma \fCenter \Delta, \lforall[x][!B(x)], !B(a)$
\RightLabel{\RightR{\lforall}}
\UnaryInf$\Gamma \fCenter \Delta, \lforall[x][!B(x)], \lforall[x][!B(x)]$
\DisplayProof
\]
Par le \olref{lem:invert-quant}, il existe une preuve~$\pi_1'$
de hauteur $\le \pheight{\pi_1}$ de
$\Gamma \fCenter \Delta, !B(c), !B(a)$
pour toute constante~$c$ absente de
$\Gamma \fCenter \Delta, \lforall[x][!B(x)], !B(a)$.
On régularise $\pi_1'$ en choisissant ses constantes
propres différentes de $a$ et de $c$.
La régularité seule ne garantirait pas cette
absence dans les inférences supérieures.
Le \olref[qua]{lem:subst-regular} s'applique
alors : $\Subst{\pi_1'}{a}{c}$ est une preuve
de $\Gamma \fCenter \Delta, !B(a), !B(a)$
de hauteur $\le \pheight{\pi_1}$.
L'hypothèse donne une preuve~$\pi_1''$
de $\Gamma \fCenter \Delta, !B(a)$.
On réapplique \RightR{\lforall} :
\[
\AxiomC{}
\RightLabel{$\pi_1''$}
\Deduce$\Gamma \fCenter \Delta, !B(a)$
\RightLabel{\RightR{\lforall}}
\UnaryInf$\Gamma \fCenter \Delta, \lforall[x][!B(x)]$
\DisplayProof
\]
La constante propre est absente de la conclusion
par la condition de l'inférence originale,
et $\pheight{\pi'} \le \pheight{\pi}$.

Supposons maintenant $!A \ident \lexists[x][!B(x)]$
principale à droite. La preuve $\pi$ a la forme
\[
\AxiomC{}
\RightLabel{$\pi_1$}
\Deduce$\Gamma \fCenter \Delta, \lexists[x][!B(x)], \lexists[x][!B(x)], !B(t)$
\RightLabel{\RightR{\lexists}}
\UnaryInf$\Gamma \fCenter \Delta, \lexists[x][!B(x)], \lexists[x][!B(x)]$
\DisplayProof
\]
L'hypothèse s'applique à~$\pi_1$ et fournit
une preuve~$\pi_1'$ de
$\Gamma \Sequent \Delta, \lexists[x][!B(x)], !B(t)$.
Aucun lemme d'inversion n'est nécessaire :
la règle a conservé les deux occurrences
de l'existentielle dans sa prémisse.
La preuve~$\pi'$ est
\[
\AxiomC{}
\RightLabel{$\pi_1'$}
\Deduce$\Gamma \fCenter \Delta,  \lexists[x][!B(x)], !B(t)$
\RightLabel{\RightR{\lexists}}
\UnaryInf$\Gamma \fCenter \Delta, \lexists[x][!B(x)]$
\DisplayProof
\]
avec $\pheight{\pi'} \le \pheight{\pi}$.
\end{proof}

\begin{prob}
Esquisser la preuve de la \olref{prop:G3c-cont-adm}
pour \LeftR{\Contraction}, en décrivant les cas
$!A \ident (!B \land !C)$ et $!A \ident (!B \lif !C)$.
\end{prob}

\begin{prob}
Esquisser la preuve de la \olref{prop:G3c-cont-adm}
pour \LeftR{\Contraction} dans les cas quantifiés
restants : $!A \ident \lforall[x][!B(x)]$
et $!A \ident \lexists[x][!B(x)]$.
\end{prob}

\end{document}
